% BHU PYQ -- Manually transcribed & error-corrected
% Paper: BSC-05A: Elements of Earth Science (Ancillary)
% Exam: B.Sc. (Hons.) II Semester Examination, 2013-14

\documentclass[11pt,a4paper]{article}
\usepackage[a4paper,top=2.5cm,bottom=2.5cm,left=2.8cm,right=2.8cm,headheight=14pt]{geometry}
\usepackage{amsmath,amssymb}
\usepackage[T1]{fontenc}
\usepackage[utf8]{inputenc}
\usepackage{lmodern,microtype,fancyhdr,enumitem,hyperref,lastpage,parskip}

\hypersetup{colorlinks=true,urlcolor=black,linkcolor=black,
  pdftitle={BSC-05A Elements of Earth Science (Ancillary) -- BHU B.Sc. Sem II 2013-14}}
\pagestyle{fancy}\fancyhf{}
\renewcommand{\headrulewidth}{0.4pt}\renewcommand{\footrulewidth}{0.4pt}
\fancyhead[L]{\small   \textit{Banaras Hindu University}}
\fancyhead[R]{\small   \textit{BSC-05A: Elements of Earth Science (Ancillary)}}
\fancyfoot[C]{\small Page  \thepage\ of \pageref{LastPage}}


\newlist{parts}{enumerate}{2}
\setlist[parts,1]{label=(\alph*),leftmargin=*,itemsep=5pt,topsep=3pt}

\newlist{romanparts}{enumerate}{2}
\setlist[romanparts,1]{label=(\roman*),leftmargin=*,itemsep=5pt,topsep=3pt}

\newcommand{\pts}[1]{\hfill   \textit{[#1]}}


\begin{document}

\begin{center}
  {\large\bfseries BANARAS HINDU UNIVERSITY}\\[2pt]
  {
\normalsize B.Sc. (Hons.) --- II Semester Examination, 2013-14}\\[6pt]
  \rule{0.8\linewidth}{0.5pt}\\[6pt]
  {\large\bfseries Paper: BSC-05A --- Elements of Earth Science (Ancillary)}\\[4pt]
  {
\normalsize Time: 3 Hours \hfill Maximum Marks: 70}\\[2pt]
  {\small Ref. No.: BS/Sem II/17 (Old)}
\end{center}
\rule{\linewidth}{0.4pt}
\smallskip



\noindent   \textit{Instructions: Answer   \textbf{five} questions in all, selecting at least   \textbf{two} from each of Sections A and B. Question No. 8 is compulsory. All questions carry equal marks.}
\bigskip

\bigskip


\noindent {\large\bfseries SECTION --- A}
\rule{\linewidth}{0.4pt}\smallskip




\noindent   \textbf{Question 1.} Describe the geological work of wind. \pts{14}

\medskip


\noindent   \textbf{Question 2.} Explain what you understand by Earth System Science and mention the various branches involved in it. \pts{14}

\medskip


\noindent   \textbf{Question 3.} What is a fault? With neat diagrams, describe the different types of faults. \pts{14}

\medskip


\noindent   \textbf{Question 4.} Write explanatory notes on   \textbf{any two} of the following: \pts{$2  \times 7 = 14$}

\begin{parts}
  \item Isostasy
  \item Mountain building
  \item Sedimentary rocks
\end{parts}

\bigskip


\noindent {\large\bfseries SECTION --- B}
\rule{\linewidth}{0.4pt}\smallskip




\noindent   \textbf{Question 5.} Describe   \textbf{any two} of the following: \pts{$2  \times 7 = 14$}

\begin{parts}
  \item Solar System
  \item Earth's magnetism
  \item Radioactivity in the Earth
\end{parts}

\medskip


\noindent   \textbf{Question 6.} Write short notes on the following: \pts{$3.5  \times 4 = 14$}

\begin{parts}
  \item Folds
  \item Moraines
  \item Deltas
  \item Types of plates
\end{parts}

\medskip


\noindent   \textbf{Question 7.} Write the chemical composition and physical properties of   \textbf{any two} of the following minerals: \pts{$2  \times 7 = 14$}

\begin{parts}
  \item Feldspar
  \item Quartz
  \item Biotite
  \item Garnet
\end{parts}

\bigskip


\noindent {\large\bfseries SECTION --- C (Compulsory)}
\rule{\linewidth}{0.4pt}\smallskip




\noindent   \textbf{Question 8.} Answer the following questions:

\begin{romanparts}
  \item   \textbf{Write True or False:} \pts{$7  \times 1 = 7$}
  
\begin{parts}
    \item Himalayas are a resultant of plate tectonics.
    \item The study of earthquakes is called Seismology.
    \item Density of the crust is more than the density of the mantle.
    \item Conglomerate is an example of an igneous rock.
    \item Diamond is the hardest of all minerals.
    \item The science of fossils is known as Paleontology.
    \item Gondwanaland is the name of a supercontinent.
  \end{parts}
  
  \item   \textbf{Fill in the blanks:} \pts{$7  \times 1 = 7$}
  
\begin{parts}
    \item V-shaped valleys are formed due to the geological work of \underline{\hspace{4cm}}.
    \item Hardness of talc on Mohs' scale is \underline{\hspace{3cm}}.
    \item Give any example of a glacial deposit: \underline{\hspace{4cm}}.
    \item The golden age of dinosaurs is the \underline{\hspace{3cm}} period in the geological time scale.
    \item The surface of erosion and non-deposition is termed as \underline{\hspace{4cm}}.
    \item Life appeared on the Earth about \underline{\hspace{3cm}} years ago.
    \item The chemical composition of calcite is \underline{\hspace{4cm}}.
  \end{parts}
\end{romanparts}

\vfill\rule{\linewidth}{0.4pt}

\begin{center}\small
\href{https://bhu-exam.vercel.app/}{Click here to visit our website}\\[4pt]
\href{https://me-aryan.vercel.app/}{Click here to visit our contributor}\\[4pt]
\href{https://quashan-me.vercel.app/}{Click here to visit our contributor}
\end{center}
\end{document}
